\section{Introduction}

Joint-embedding predictive architectures (JEPAs) learn by predicting the representation of
a masked target region from the representation of a visible context region
\citep{assran2023ijepa}. Because both sides of the objective are produced by trainable
networks, a constant encoder satisfies the objective perfectly while carrying no
information about the input. Preventing this degenerate solution is a design requirement
rather than an incidental concern.

Published JEPAs prevent collapse through a collection of empirical mechanisms:
stop-gradient on the target branch, an exponential-moving-average (EMA) target encoder,
masking design, and normalization choices \citep{grill2020byol,chen2021simsiam}. These work
in practice but come with no guarantee. LeJEPA \citep{lejepa2025} proposes SIGReg, a
distributional regularizer intended to replace them.

This report asks a narrow question about that claim. If SIGReg removes the need for
stop-gradient, the two should remain interchangeable no matter how long training runs. We
test whether that holds as training duration grows, and, separately, whether the cheap
diagnostics normally used to detect collapse give any warning before downstream probe
accuracy degrades.
